\documentclass[11pt,letterpaper]{article}

% ── Geometry ──────────────────────────────────────────────────────────────────
\usepackage[margin=1in]{geometry}

% ── Pagination ────────────────────────────────────────────────────────────────
\usepackage[all]{nowidow}    % Require ≥2 lines on each side of a page break
\brokenpenalty=10000         % No hyphenated words across page breaks
\raggedbottom                % Allow uneven page bottoms rather than stretching

% ── Fonts ─────────────────────────────────────────────────────────────────────
\usepackage[T1]{fontenc}
\usepackage{newpxtext}     % Palatino-based serif (TeX Gyre Pagella)
\usepackage{newpxmath}     % Matching math font

% ── Colors ────────────────────────────────────────────────────────────────────
\usepackage[dvipsnames]{xcolor}
\definecolor{SectionBlue}{HTML}{1B3A5C}
\definecolor{AccentGray}{HTML}{4A4A4A}
\definecolor{QuoteBorder}{HTML}{8B9DAF}
\definecolor{RiskRed}{HTML}{C0392B}
\definecolor{RiskAmber}{HTML}{E67E22}
\definecolor{RiskGreen}{HTML}{27AE60}
\definecolor{BoxBg}{HTML}{F7F8FA}

% ── Tables ────────────────────────────────────────────────────────────────────
\usepackage{booktabs}
\usepackage{tabularx}

% ── Boxes ─────────────────────────────────────────────────────────────────────
\usepackage[framemethod=tikz]{mdframed}

% ── Lists ─────────────────────────────────────────────────────────────────────
\usepackage{enumitem}
\usepackage{needspace}

% ── Hyperlinks ────────────────────────────────────────────────────────────────
\usepackage{hyperref}
\usepackage{xurl}  % break bibliography URLs at any character, not just at slashes
\hypersetup{
  colorlinks=true,
  linkcolor=SectionBlue,
  urlcolor=SectionBlue,
  citecolor=SectionBlue,
  pdfauthor={JM Freeman},
  pdftitle={Quarry Answer Synthesis: PR/FAQ},
  pdfsubject={Working Backwards: Product Discovery},
}

% ── Bibliography ─────────────────────────────────────────────────────────────
\usepackage[style=numeric,backend=biber,sorting=none]{biblatex}
\addbibresource{prfaq-answer-synthesis.bib}

% ── Section styling ───────────────────────────────────────────────────────────
\usepackage{titlesec}
\titleformat{\section}
  {\needspace{6\baselineskip}\Large\bfseries\color{SectionBlue}}
  {}
  {0pt}
  {}

\titleformat{\subsection}
  {\needspace{5\baselineskip}\large\bfseries\color{AccentGray}}
  {}
  {0pt}
  {}

\titleformat{\subsubsection}
  {\normalsize\bfseries\color{AccentGray}}
  {}
  {0pt}
  {}
\titlespacing*{\subsubsection}{0pt}{1em}{0.5em}[0pt]

% ── Document stage ───────────────────────────────────────────────────────────
% Stage calibrates evidence expectations across discovery, drafting, peer review, and meetings.
% Valid values: hypothesis | validated | growth
\newcommand{\prfaqstage}[1]{\def\theprfaqstage{#1}}
\prfaqstage{hypothesis}  % Default: pre-validation

% ── Document version ─────────────────────────────────────────────────────────
% Tracks iteration count. Major for structural changes, minor for feedback edits.
% /prfaq:feedback auto-increments after each application.
\newcommand{\prfaqversion}[2]{\def\theprfaqmajor{#1}\def\theprfaqminor{#2}}
\prfaqversion{1}{1}

% ── Header/footer ────────────────────────────────────────────────────────────
\usepackage{fancyhdr}
\pagestyle{fancy}
\fancyhf{}
\fancyhead[L]{\footnotesize\textcolor{AccentGray}{Stage: \theprfaqstage{} \enspace\textbar\enspace v\theprfaqmajor.\theprfaqminor}}
\fancyhead[R]{\footnotesize\textcolor{AccentGray}{\thepage}}
\renewcommand{\headrulewidth}{0pt}

% ══════════════════════════════════════════════════════════════════════════════
% Custom environments
% ══════════════════════════════════════════════════════════════════════════════

% ── \prsection{title} — styled press release section header ──────────────────
\newcommand{\prsection}[1]{%
  \vspace{1em}%
  {\large\bfseries\color{SectionBlue}#1}%
  \vspace{0.3em}\par%
}

% ── customerquote — indented, italic customer quote ──────────────────────────
\newmdenv[
  topline=false,
  bottomline=false,
  rightline=false,
  linewidth=3pt,
  linecolor=QuoteBorder,
  backgroundcolor=BoxBg,
  innerleftmargin=12pt,
  innerrightmargin=12pt,
  innertopmargin=8pt,
  innerbottommargin=8pt,
  skipabove=\baselineskip,
  skipbelow=\baselineskip,
]{customerquote}

% ── spokespersonquote — indented spokesperson quote ──────────────────────────
\newmdenv[
  topline=false,
  bottomline=false,
  rightline=false,
  linewidth=3pt,
  linecolor=SectionBlue,
  backgroundcolor=BoxBg,
  innerleftmargin=12pt,
  innerrightmargin=12pt,
  innertopmargin=8pt,
  innerbottommargin=8pt,
  skipabove=\baselineskip,
  skipbelow=\baselineskip,
]{spokespersonquote}

% ── faqpair — numbered Q&A pair with bold question ─────────────────────────
% Label FAQ pairs that explain judgment calls: \begin{faqpair}{Question}\label{faq:slug}
% Then cross-reference from the press release: (see \faqref{faq:slug})
\newcounter{faqnum}
\newenvironment{faqpair}[1]{%
  \needspace{4\baselineskip}%
  \vspace{0.5em}%
  \refstepcounter{faqnum}%
  \noindent\textbf{\color{SectionBlue}Q\thefaqnum: #1}\par\nopagebreak[4]%
  \vspace{0.2em}%
  \begin{adjustwidth}{1em}{0pt}%
  \setlength{\parindent}{0pt}%
}{%
  \end{adjustwidth}%
  \vspace{0.5em}%
}
% ── \faqref{faq:slug} — clickable "FAQ 3" reference ──────────────────────────
\newcommand{\faqref}[1]{\hyperref[#1]{FAQ~\ref*{#1}}}

% ── \featureitem — numbered, referenceable feature entry ──────────────────────
% Use inside enumerate: \featureitem{Name}{Rationale}\label{feat:slug}
% Then cross-reference: (see \featureref{feat:slug})
\newcounter{featurenum}
\newcommand{\featureitem}[2]{%
  \refstepcounter{featurenum}%
  \item[\textbf{\color{SectionBlue}F\thefeaturenum.}]%
  \textbf{#1}: #2%
}
\newcommand{\featureref}[1]{\hyperref[#1]{Feature~\ref*{#1}}}

% ── adjustwidth (for faqpair indentation) ────────────────────────────────────
\usepackage{changepage}

% ══════════════════════════════════════════════════════════════════════════════
% Document
% ══════════════════════════════════════════════════════════════════════════════

\begin{document}

\begin{center}
  {\LARGE\bfseries\color{SectionBlue} Quarry Now Answers the Question Instead of Listing Ten Passages} \\[0.8em]
  {\large\color{AccentGray} quarryd turns its own search results into a short answer with numbered citations, using any OpenAI-compatible model on the same Mac or in the cloud} \\[0.3em]
  {\large\color{AccentGray} The menu bar app and the \texttt{quarry} CLI show the answer in one keystroke, against a local or a remote daemon}
\end{center}

\vspace{0.5em}

% ══════════════════════════════════════════════════════════════════════════════
% PRESS RELEASE
% ══════════════════════════════════════════════════════════════════════════════

\section*{Press Release}

PORTLAND, OR, April 2027. Punt Labs today added answer synthesis to Quarry,
its local semantic search engine. Developers who search their own Quarry index
from the macOS menu bar or the terminal can now press one key and get a three
to six sentence answer, written from the passages Quarry already found, with a
numbered citation on every claim. Each citation opens the source page. The
answer is written by whichever OpenAI-compatible model the user points
\texttt{quarryd} at: Ollama or LM Studio on the same machine for fully local
use, or a hosted Haiku-class model. Chat assistants answer from documents
the user uploads to them; Quarry answers from the index it already keeps, and
search stays as fast as it was (see \faqref{faq:vs-chat}).

\prsection{Problem}

A Quarry search returns up to ten passages, and each passage can be up to
1{,}800 characters long. A full page of results is about 18{,}000 characters,
roughly the length of a ten-page memo. The menu bar app shows the first 200
characters of each result. To answer a question like ``which port does the
staging daemon listen on, and who changed it?'', the user opens results one at
a time, reads each passage, and assembles the answer in their head.

Inside Claude Code this is not a problem, because the agent reads the
passages for the user through Quarry's MCP tools. Outside an agent session it
is: a quick lookup from the menu bar or a shell ends in reading. The usual
workaround is to copy passages into a chat window, which loses the link back
to the source and sends the text to whichever service the chat window uses
(see \faqref{faq:evidence}).

\prsection{Solution}

The user types a question into the menu bar search panel as they do today.
Results appear as they type, with no model call. When they want an answer, they
press Command-Return. \texttt{quarryd} runs the same search, sends the top
passages to the configured model with instructions to answer only from those
passages and to cite them by number, and returns the answer together with the
results it used. The answer appears in a card above the result list. Clicking
\texttt{[3]} opens result 3 in the existing detail view, so every sentence is
one click from its source. If the passages do not cover the question, the
answer says so instead of guessing (see \faqref{faq:wrong}).

The same answer is available from the terminal with
\texttt{quarry answer}. Because the daemon does the work, the menu bar
app gets answers whether it is connected to the local \texttt{quarryd} or to a
remote one, and the app never holds a model key (see \faqref{faq:why-daemon}).

\prsection{Customer Quote}

\begin{customerquote}
\textit{``I keep runbooks, design records and two years of incident notes in
Quarry. A search used to give me ten snippets, and I'd open four of them to
piece the answer together. Now I press Command-Return and get two sentences
with a [2] and a [5] after them. I still click [5] before I change anything in
production, and it takes me straight to the paragraph. The model runs in Ollama
on my laptop, so the incident notes stay on my laptop too.''}

\raggedleft \textbf{Priya Raman}, Platform Engineer at a 40-person logistics software company
\end{customerquote}

\prsection{Getting Started}

Users who already run Quarry need about five minutes:

\begin{enumerate}[leftmargin=2em]
  \item \textbf{Pick a model.} Run one locally with Ollama, LM Studio or
        llama.cpp, or use a hosted OpenAI-compatible endpoint and its API
        key.
  \item \textbf{Point \texttt{quarryd} at it.} Add the endpoint URL and model
        name to Quarry's configuration and restart the daemon. Synthesis stays
        off until this step, and nothing else changes.
  \item \textbf{Ask.} Press Command-Return in the menu bar search panel, or run
        \texttt{quarry answer "your question"}.
\end{enumerate}

\prsection{Spokesperson Quote}

\begin{spokespersonquote}
``Quarry already finds the right passages. The step it left to the user was
reading ten of them to get one fact. That step is a small, well-bounded job for
a small model: answer from these passages and say which ones. We put it in the
daemon so every Quarry client gets the same answer from the same configuration,
and we made the model a setting so users choose whether their notes ever leave
the machine. Search stays exactly as fast as it was; the answer is one
keystroke on top.''

\raggedleft \textbf{JM Freeman}, Creator of Quarry
\end{spokespersonquote}

\prsection{Call to Action}

Answer synthesis ships in the next minor releases of \texttt{punt-quarry} and
Quarry Menu Bar, free and MIT-licensed like the rest of Quarry. The model is
the user's choice and the user's cost: nothing for a local model, and about
half a US cent per answer for Claude Haiku 4.5 at list price
(see \faqref{faq:cost}). Install or upgrade with
\texttt{uv tool install punt-quarry}.

% ══════════════════════════════════════════════════════════════════════════════
% FAQ
% ══════════════════════════════════════════════════════════════════════════════

\newpage
\section*{Frequently Asked Questions}

\subsection*{External FAQs}

\begin{faqpair}{What is answer synthesis and who is it for?}\label{faq:what}
  It is a short written answer to a Quarry search, with numbered citations
  back to the results. It is for people who already keep their documents,
  notes and captured research in Quarry and look things up from the menu bar
  app or the terminal, outside an agent session. The job it does is the last
  step of a lookup: turning ten relevant passages into the one or two facts the
  user wanted.
\end{faqpair}

\begin{faqpair}{How is this different from pasting the results into ChatGPT or Claude?}\label{faq:vs-chat}
  The result is similar; the work is gone. There is no copying, the citations
  stay linked to the source pages, and the user decides once, in
  configuration, which model sees their passages. Hosted notebook products
  such as NotebookLM also answer with citations~\cite{google2026notebooklm},
  but they answer from documents uploaded to that service. Quarry answers from
  the local index that already holds the user's files, web captures and session
  notes.
\end{faqpair}

\begin{faqpair}{I use Quarry through Claude Code. Does this change anything for me?}\label{faq:claude-code}
  No. In Claude Code the agent already reads Quarry's results and writes the
  answer itself. Quarry's MCP tools keep returning passages, because a
  pre-written answer would cost the agent tokens and hide the source text it
  reasons over (see \featureref{feat:no-mcp}).
\end{faqpair}

\begin{faqpair}{How do I get started?}\label{faq:start}
  Pick a model, add its endpoint URL and model name to Quarry's configuration,
  restart \texttt{quarryd}, then press Command-Return in the menu bar search
  panel or run \texttt{quarry answer}. Until a model is configured, the
  menu bar app hides the Answer action and the CLI command reports that synthesis
  is not configured.
\end{faqpair}

\begin{faqpair}{How much does it cost?}\label{faq:cost}
  Quarry is free. The model costs whatever the chosen endpoint charges. A local
  model costs nothing beyond electricity. For a hosted model, one answer sends
  about eight passages plus instructions, roughly 4{,}000 input tokens at four
  characters per token, and returns about 250 output tokens. At Claude Haiku
  4.5 list prices of \$1 per million input tokens and \$5 per million output
  tokens~\cite{anthropic2026pricing}, that is about \$0.005 per answer, or
  \$5 for a thousand answers.
\end{faqpair}

\begin{faqpair}{What happens to my data?}\label{faq:privacy}
  With a local model, nothing leaves the machine. With a hosted model, each
  answer sends that question and the top passages for it to the configured
  endpoint, and nothing else: the index, other documents and the search history
  stay local. Every answer response names the endpoint host that wrote it, and
  the menu bar card shows that host under the answer. Quarry stores no answers
  and sends nothing to Punt Labs. Session captures are passed through Quarry's
  secret scrubber when they are ingested, but the rest of the user's corpus is
  sent as indexed, so users with sensitive files should choose a local model.
\end{faqpair}

\begin{faqpair}{What if the answer is wrong?}\label{faq:wrong}
  The model is told to answer only from the numbered passages, to cite each
  claim, and to say when the passages do not answer the question. Quarry then
  checks the citations: a number that does not match a returned result is
  removed before the answer reaches the user. The results are always shown
  under the answer, so the user can check any sentence in one click. A small
  model can still misread a passage; the citations make that visible rather
  than preventing it (see \faqref{faq:tech-risks}).
\end{faqpair}

\begin{faqpair}{I connect the menu bar app to Quarry on another machine. Whose model answers?}\label{faq:remote}
  The remote daemon's. The model call runs on the machine that runs
  \texttt{quarryd}, with that machine's configuration and key. The menu bar app
  sends the question with the same bearer token it already uses for search and
  stores no model settings of its own. If the remote daemon has no model
  configured, the Answer action is hidden.
\end{faqpair}

\begin{faqpair}{Will this make search slower?}\label{faq:speed}
  No. Search-as-you-type keeps calling the existing search endpoint with no
  model in the path. An answer is a separate request that runs only when the
  user asks for one, and if the model is slow or down, the answer card shows
  the error while search results keep working.
\end{faqpair}

\subsection*{Internal FAQs}

\subsubsection*{Value \& Market}

\begin{faqpair}{What is the total addressable market?}\label{faq:tam}
  Quarry is free and open source, so the relevant market is Quarry's own users,
  and the feature is judged by whether it makes Quarry more useful to them.
  \texttt{punt-quarry} had 469 PyPI downloads in the month to
  2026-10-25~\cite{pypistats2026quarry}; downloads include CI and repeat
  installs, so active users are fewer. The menu bar app is macOS-only and has
  no download or usage counter, so its share of that population is unknown.
  The upper bound on reach is every Quarry user who searches outside an agent
  session, through the menu bar app or the CLI.
\end{faqpair}

\begin{faqpair}{What evidence do we have that users want this?}\label{faq:evidence}
  Little beyond the author's own use. The request comes from the Quarry author
  using the menu bar app daily, which is the weakest kind of evidence: the
  customer is the team. Two indirect signals support it. First, the size of the
  results: ten passages of up to 1{,}800 characters, of which the menu bar
  shows 200, leave reading as the main cost of a lookup. Second, products built
  on the same pattern, such as NotebookLM~\cite{google2026notebooklm}, answer
  from a user's sources with citations. Neither shows that Quarry users want an
  answer more than they want the passages.
\end{faqpair}

\begin{faqpair}{Who are the alternatives and why would a user pick this?}\label{faq:competitors}
  The real alternatives are habits. (1) Reading the results, which costs
  nothing and needs no model. (2) Asking an agent that has Quarry's MCP tools,
  which already synthesizes but means opening a session for a quick lookup.
  (3) Pasting passages into a chat assistant, which loses citations and sends
  text to that service. (4) Hosted notebook products~\cite{google2026notebooklm},
  which require uploading documents and do not see Quarry's captures. Quarry's
  advantage is position: it already holds the corpus and the ranked passages,
  so the answer is one model call away. Any of these alternatives could add
  something similar; none of them has the user's local index.
\end{faqpair}

\begin{faqpair}{What is the next step to validate the vision?}\label{faq:next-step}
  Before building the menu bar card, ship the daemon route and the CLI command,
  then run the eval spike in \faqref{faq:tech-risks} and use
  \texttt{quarry answer} for two weeks of normal work. Record how often
  an answer is requested compared with a plain search, how often the answer
  says the passages do not cover the question, and how often the user still
  opens a result. If the author stops using \texttt{quarry answer} within the two
  weeks, the menu bar work does not start.
\end{faqpair}

\subsubsection*{Technical}

\begin{faqpair}{Why put synthesis in \texttt{quarryd} instead of the menu bar app?}\label{faq:why-daemon}
  Three reasons. Every client gets the same behavior: the menu bar app, the CLI
  and any future client call one route, and a remote menu bar connection gets
  answers without its own model setup. Keys live in one place, on the machine
  that runs the daemon; the menu bar app has no settings screen today and would
  otherwise need one plus Keychain storage. And the daemon already has the
  passages in memory when it searches, so it can send them to the model without
  a second round trip. The cost is that the daemon now makes outbound HTTP calls
  and holds a model key, and a remote daemon spends its owner's tokens on
  behalf of every caller holding its bearer token.
\end{faqpair}

\begin{faqpair}{Why not run a generative model inside \texttt{quarryd}, like the embedding model?}\label{faq:why-not-bundled}
  The embedding model is a small encoder that ONNX Runtime runs on the CPU
  inside the daemon. A generative model of useful quality is many times larger,
  would add gigabytes of download and resident memory to a daemon that is
  always running, and on macOS would run on the CPU because Quarry's ONNX path
  has no GPU provider there. Generation also needs chat templates, a sampling
  loop and stop handling that Quarry would have to own. Ollama, LM Studio and
  llama.cpp already serve local models behind an OpenAI-compatible
  API~\cite{ollama2026openai,lmstudio2026openai,llamacpp2026server} and use the
  Mac's GPU, so one HTTP client in Quarry gets local synthesis without Quarry
  managing a model.
\end{faqpair}

\begin{faqpair}{What are the major technical risks?}\label{faq:tech-risks}
  \textbf{Answer quality from small local models (high, unknown).} A 1--4B
  parameter model may cite the wrong passage or merge two passages into a claim
  neither makes. Mitigation: a spike that runs Quarry's eval fixtures through
  two local models and one hosted model, scored for citation correctness by
  hand, before choosing a recommended default. \textbf{Latency on a laptop
  (medium, unknown).} A local model's time to a full answer on an M-series Mac
  is not measured; the spike records it. \textbf{Provider compatibility
  (medium).} ``OpenAI-compatible'' varies between servers. Anthropic describes
  its compatibility layer as intended for testing and not a long-term or
  production-ready solution for most use cases~\cite{anthropic2026openaisdk},
  so the client must stay within the core chat completions fields and treat
  provider differences as test cases. \textbf{Key delivery to a background
  daemon (medium).} \texttt{quarryd} runs as a service, so keys exported in a
  shell do not reach it; the design must name where the key comes from on macOS
  and Linux. \textbf{Cost exposure on a shared daemon (low for personal use).}
  Any holder of a remote daemon's bearer token can spend its model budget; a
  per-daemon request limit bounds this.
\end{faqpair}

\begin{faqpair}{What dependencies exist on other systems?}\label{faq:dependencies}
  The OpenAI chat completions request and response format, as implemented by
  the endpoint the user chooses. \texttt{httpx}, already a Quarry dependency,
  for the outbound call. The existing \texttt{SearchService}, authorization and
  collection filtering inside \texttt{quarryd}, reused unchanged. On the client
  side, the menu bar app's existing \texttt{QuarryClient} and detail view. No
  new third-party service is required: a user with no model configured sees no
  change.
\end{faqpair}

\begin{faqpair}{What is the delivery order?}\label{faq:timeline}
  (1) Daemon: configuration, the OpenAI-compatible client, the answer route
  with citation checking, the OpenAPI update and tests against a fake endpoint.
  (2) CLI: \texttt{quarry answer}. (3) Eval spike and two weeks of use
  (\faqref{faq:next-step}); this gates the next phase. (4) Menu bar: the Answer
  action, the answer card with clickable citations, and the endpoint host
  label. (5) Streaming, if the measured latency calls for it. If scope
  compresses, streaming is cut first, then the CLI command; the daemon route and
  the menu bar card are the feature.
\end{faqpair}

\begin{faqpair}{What breaks first if it is used heavily?}\label{faq:scaling}
  The model endpoint. Each answer is one search plus one model call, so the
  daemon's own load barely changes, but a local model serves one request at a
  time and a hosted one has rate limits. The daemon needs a timeout on the model
  call, a cap on concurrent answer requests, and an error the client can show
  without blocking search.
\end{faqpair}

\subsubsection*{Business}

\begin{faqpair}{What is the revenue model?}\label{faq:revenue}
  None. Quarry is free and MIT-licensed. The value to Punt Labs is a menu bar
  app that is useful on its own for quick lookups, which brings users to the
  rest of Quarry. Punt Labs pays no model costs: the user's endpoint and key
  are used for every call.
\end{faqpair}

\begin{faqpair}{How will we measure success?}\label{faq:metrics}
  Quarry's local recall telemetry already records searches; answer requests
  would be recorded the same way, on the user's machine only. (1) Answer
  requests as a share of searches, per active user per week; below 10\% after
  the two-week trial suggests users prefer passages. (2) Share of answers that
  say the passages do not cover the question; above 30\% points at retrieval,
  not synthesis. (3) Citation clicks per answer, a sign that users check
  answers. (4) Median time to a full answer for local and hosted models; above
  five seconds makes streaming a must. (5) Model call error rate.
\end{faqpair}

\begin{faqpair}{Why now?}\label{faq:why-now}
  Three things are in place. The common local model servers expose the same
  OpenAI-compatible
  API~\cite{ollama2026openai,lmstudio2026openai,llamacpp2026server}, so one
  client covers local and hosted models. Haiku-class models answer a
  4{,}000-token question for about half a cent~\cite{anthropic2026pricing}. And
  Quarry is now daemon-first, with every client going through \texttt{quarryd},
  so a feature added there reaches the menu bar app, the CLI and remote
  connections at once.
\end{faqpair}

\begin{faqpair}{What are we not building?}\label{faq:not-building}
  Multi-turn chat, a synthesis tool in MCP, a model bundled inside Quarry,
  model settings in the menu bar app, and answers that run on every keystroke.
  The Feature Appendix lists each with its reason.
\end{faqpair}

% ══════════════════════════════════════════════════════════════════════════════
% FOUR RISKS ASSESSMENT
% ══════════════════════════════════════════════════════════════════════════════

\newpage
\section*{Risk Assessment}

\renewcommand{\arraystretch}{1.6}
\begin{tabularx}{\textwidth}{@{} l l X @{}}
\toprule
\textbf{\color{SectionBlue}Risk} & \textbf{\color{SectionBlue}Rating} & \textbf{\color{SectionBlue}Assessment} \\
\midrule
Value       & \textcolor{RiskAmber}{Medium} & Demand evidence is the author's own use (\faqref{faq:evidence}). The two-week \texttt{quarry answer} trial is the test, and it runs before the menu bar work. \\
Usability   & \textcolor{RiskAmber}{Medium} & One keystroke to use, but setup asks the user to choose and run a model and edit configuration. Users without a model see no change. \\
Feasibility & \textcolor{RiskAmber}{Medium} & The plumbing is one route and one HTTP call over existing search. Answer quality and latency from small local models are unmeasured (\faqref{faq:tech-risks}). \\
Viability   & \textcolor{RiskGreen}{Low}    & No cost to Punt Labs; users pay their own model costs. Main exposure is a shared remote daemon spending its owner's tokens. \\
\bottomrule
\end{tabularx}
\renewcommand{\arraystretch}{1.0}

% ══════════════════════════════════════════════════════════════════════════════
% FEATURE APPENDIX
% ══════════════════════════════════════════════════════════════════════════════

\newpage
\section*{Feature Appendix}

\subsection*{Must Do}

\begin{enumerate}[nosep,leftmargin=2.5em]
  \featureitem{Answer route in \texttt{quarryd}}{a \texttt{/v1} route separate from search that reuses \texttt{SearchService}, authorization and collection filtering, and returns the answer with the results it cites}\label{feat:route}
  \featureitem{OpenAI-compatible client}{one chat completions call over \texttt{httpx} with a timeout and a concurrency cap; works with Ollama, LM Studio, llama.cpp and hosted endpoints}\label{feat:client}
  \featureitem{Off until configured}{endpoint URL, model and key come from daemon configuration; with none set, the route reports ``not configured'' and the daemon's status tells clients to hide the feature}\label{feat:config}
  \featureitem{Grounded prompt and citation check}{numbered passages, answer-only-from-sources instructions, an explicit ``not covered'' reply, and server-side removal of citation numbers that match no result}\label{feat:grounding}
  \featureitem{Endpoint disclosure}{every answer names the endpoint host that wrote it, and clients display it}\label{feat:disclosure}
  \featureitem{Menu bar Answer card}{an explicit Command-Return or button trigger, a card above the results, clickable citations that open the existing detail view, and cancellation when the query changes}\label{feat:menubar}
\end{enumerate}

\subsection*{Should Do}

\begin{enumerate}[nosep,leftmargin=2.5em]
  \featureitem{\texttt{quarry answer}}{the CLI surface, cheap once the route exists, and the vehicle for the validation trial}\label{feat:cli}
  \featureitem{Answer quality eval}{run Quarry's eval fixtures through candidate models and score citation correctness, to choose a recommended default}\label{feat:eval}
  \featureitem{Platform secret lookup}{read the model key from the macOS Keychain or \texttt{pass} so it need not sit in a file the daemon reads}\label{feat:secrets}
  \featureitem{Answer telemetry}{record answer requests, ``not covered'' replies and latency in Quarry's local recall telemetry}\label{feat:telemetry}
  \featureitem{Streaming}{show the answer as it is written, if measured latency makes waiting noticeable}\label{feat:streaming}
\end{enumerate}

\subsection*{Won't Do}

\begin{enumerate}[nosep,leftmargin=2.5em]
  \featureitem{Synthesis in MCP}{the MCP caller is already a model; a pre-written answer costs it tokens and hides the passages it reasons over}\label{feat:no-mcp}
  \featureitem{A bundled generative model}{gigabytes of download and memory in an always-on daemon, CPU-only on macOS, and a sampling loop to own, when local servers already exist (\faqref{faq:why-not-bundled})}\label{feat:no-bundle}
  \featureitem{Model settings in the menu bar app}{keys and endpoints belong to the daemon that makes the call; the app stays a thin client}\label{feat:no-app-config}
  \featureitem{Answers on every keystroke}{search-as-you-type would fire a model call for every half-typed query}\label{feat:no-auto}
  \featureitem{Multi-turn chat}{follow-up questions turn a lookup tool into a chat client, a different product with its own state and UI}\label{feat:no-chat}
  \featureitem{Native, non-OpenAI provider APIs}{one request format keeps the client small; providers without an OpenAI-compatible endpoint can be reached through a gateway}\label{feat:no-native}
\end{enumerate}

% ══════════════════════════════════════════════════════════════════════════════
% REFERENCES
% ══════════════════════════════════════════════════════════════════════════════

\newpage
\printbibliography[title={References}]

\end{document}
